\scanpage{32}
% archive sheet number: 31
\begingroup
\makeatletter\tagsleft@true\makeatother

\margin{\uline{Version II}}
\noindent Décompositions $n$-aires d'un espace topologique\unsure{.}

1.1. Soit $X$ un espace topologique. On appelle \uline{décomposition binaire} de $X$ une paire (\unsure{notée|noté} indiciellement)
\[
\unsure{\pi|\Pi}=(\Sigma^\varepsilon)_{\varepsilon\in\omega}\qquad(\omega\ \text{ens à deux éléments})
\]
\unsure{dé|de|d'e} \deleted{ensembles} parties \uline{fermées} (distinctes) de $X$, de réunion $X$
\[
X=\bigcup_{\varepsilon\in\omega}\Sigma^\varepsilon.\tag*{(1)}
\]

On suppose que
\[
\Sigma^\varepsilon\ne\varnothing\qquad(\forall\varepsilon\in\omega)\tag*{(2)}
\]
ce qui signifie aussi que
\[
\pi\ne\{\varnothing,X\}\tag*{$(2')$}
\]
et implique que $X\ne\varnothing$, et même $\mathrm{card}\,X\geq 2$, puisque l'on a supposé
\[
\unsure{\deleted{X|\Sigma}}\Sigma^\varepsilon\unsure{\deleted{=|\ne}}\ne\Sigma^{\varepsilon'}\qquad\text{si }\varepsilon,\varepsilon'\in\omega,\quad\varepsilon\ne\varepsilon'.\tag*{(\unsure{3|s|$\wedge$})}
\]

On n'exclut pas le cas où on aurait une inclusion
\[
\Sigma^{\varepsilon'}\subset\Sigma^\varepsilon
\]
ce qui équivaut, par (1), à la relation
\[
\Sigma^\varepsilon=X.
\]

Les $\Sigma^\varepsilon$ $(\varepsilon\in\omega)$ s'appellent les \uline{\unsure{hémisphères|hémisphéres}} de la décomposition binaire, \deleted{la} et l'ensemble
\[
E=\bigcap_{\deleted{\omega}\varepsilon\in\omega}\Sigma^\omega\tag*{(4)}
\]
s'appelle l'\uline{équateur} de la décomposition binaire.

On n'exclut pas le cas où on aurait $E=\varnothing$, ce qui indique que les $\Sigma^\varepsilon$ \inserted{$(\varepsilon\in\omega)$} sont \inserted{complémentaires l'un de l'autre, \unsure{donc}} à la fois fermés et ouverts -- donc que $X$ est disconnexe.
\endgroup
